\documentclass[11pt]{article}

\usepackage{geometry}
\usepackage{hyperref}
\usepackage{tikz}
\usetikzlibrary{arrows.meta, positioning, calc}
\geometry{margin=1in}

\title{Emergent State Machines:\\A Claim--Evidence--Reasoning Example}
\author{}
\date{}

\begin{document}

\maketitle

\section{Purpose}

This document provides a simple example of how an Emergent State Machine (ESM) can be applied to student writing.

The example focuses on a common instructional task in middle and secondary education: helping students produce writing that includes a clear claim, supporting evidence, and explicit reasoning.

\section{The Problem}

Teachers often see the same structural problems repeatedly in student writing.

For example, a student may:

\begin{itemize}
\item make a claim without evidence
\item list evidence without explaining it
\item provide reasoning without grounding it in a claim
\end{itemize}

In many classrooms, the teacher must repeatedly diagnose these problems by hand and decide what support to give next.

An ESM can support this process by making the student's current argumentative situation explicit. See Figure~\ref{fig:cer-esm-flow}.

\input{cer_diagram}

\section{What the ESM Monitors}

In this example, the ESM monitors one narrow domain:

\begin{quote}
\textit{the structure of claim--evidence--reasoning in the student's current writing}
\end{quote}

The system is not evaluating the entire student, the entire assignment, or the entire course. It is evaluating one specific writing situation.

\section{The ESM Pipeline}

The ESM operates through the same general loop used in other domains.

\subsection*{Observations}

The system receives student writing as input.

Examples of observations include:

\begin{itemize}
\item student sentence text
\item paragraph text
\item revision attempts
\end{itemize}

\subsection*{Signals}

From those observations, the system computes direct and derived signals.

Examples include:

\begin{itemize}
\item claim present
\item evidence present
\item reasoning present
\item evidence--reasoning imbalance
\item repeated missing-claim pattern
\end{itemize}

\subsection*{State Vector}

The signals are assembled into an interpretable state representing the student's current argumentative situation.

Examples of states include:

\begin{itemize}
\item claim missing
\item claim present but unsupported
\item evidence present without reasoning
\item complete CER structure
\end{itemize}

\subsection*{Gate}

A gate determines whether certain forms of evaluation should proceed.

For example, if no claim is present, the system may defer deeper evidence or reasoning evaluation and instead activate claim-oriented support.

\subsection*{Policy Projection}

The state is evaluated relative to instructional policy dimensions.

Examples include:

\begin{itemize}
\item missing foundational structure
\item ready for evidence support
\item ready for reasoning support
\item ready for revision or enrichment
\end{itemize}

\subsection*{Policy Selection}

Based on the projected state, the system selects a next instructional action.

Examples include:

\begin{itemize}
\item prompt the student to write a claim
\item ask for supporting evidence
\item ask the student to explain how the evidence supports the claim
\item direct the student to revise a weak reasoning statement
\end{itemize}

\subsection*{State Evolution}

As the student revises, new observations arrive and the state updates.

Over time, the student's writing forms a trajectory through the state space.

\section{Playbooks}

A playbook defines how the system guides the student over a sequence of states.

For example, a basic CER playbook might guide the student through:

\begin{enumerate}
\item producing a claim
\item adding evidence
\item adding reasoning
\item strengthening reasoning
\end{enumerate}

The playbook does not replace the ESM. Instead, it uses ESM state evaluations to decide what instructional support should come next.

\section{Breadcrumbs}

When the student reaches an appropriate transition point, the system may present breadcrumbs representing possible next steps.

Examples might include:

\begin{itemize}
\item Revise your reasoning
\item Try a paragraph-level argument
\item Apply this structure to a science explanation
\end{itemize}

Breadcrumbs make the next transition visible while preserving structured guidance.

\section{Why This Matters}

The value of the ESM in this context is not that it ``grades'' writing automatically.

Its value is that it can:

\begin{itemize}
\item make structural writing situations explicit
\item reduce repetitive diagnostic work for teachers
\item support students with immediate, targeted feedback
\item guide students through interpretable writing states
\end{itemize}

This allows teachers to spend more time on higher-level instructional work such as interpretation, discussion, extension, and creative development.

\section{Summary}

In this example, the ESM functions as a deterministic, inspectable reasoning structure for argumentative writing.

It transforms writing observations into signals, signals into state, and state into policy-guided instructional action.

The result is a system that supports writing development through explicit situational reasoning rather than opaque scoring.

\end{document}

